\newpage
\section{Future Work}

A clear next step is improving dataset coverage and cross-camera generalization. 
The current detector was trained on imagery from only 12 intersections, and the 
cross-camera benchmark showed that performance dropped on unseen views. Expanding 
the dataset to include a wider range of camera locations, boroughs, lighting 
conditions, seasons, and viewing angles would make the system more robust for 
citywide deployment. Additional targeted annotation would also be especially 
valuable for underrepresented classes such as \verb|motorcycle|, which had very 
few training examples and remained the least reliable class.

A second priority is improving how frame-level detections are converted into traffic 
counts. The ATR comparison showed that the current snapshot-based approach was better
at recovering broad hourly traffic patterns than exact count magnitude. Future 
systems should therefore move beyond independent still-frame detection and incorporate 
temporal modeling. In particular, adding a lightweight multi-object tracking method 
would help prevent repeated counting of persistent vehicles and reduce missed counts 
from vehicles that pass through the scene between snapshots. A statistical calibration
step that maps image-derived counts to ATR-style flow measurements would also help 
bridge the gap between detector output and true traffic volume.

A final direction is strengthening the connection between vehicle detection and 
emissions estimation. Although the current 6-class taxonomy was designed to align 
with MOVES, full MOVES-based emissions modeling also depends on operational variables
such as vehicle speed and vehicle miles traveled. Integrating camera-based detections
with external traffic data, roadway geometry, or speed estimates would make the 
pipeline more useful for downstream environmental analysis. Additional extensions 
could also explore predicting semantic attributes such as taxi or delivery status 
while preserving the core MOVES-aligned vehicle classes. Together, these improvements
would move the pipeline from snapshot-based fleet characterization toward a more 
complete system for intersection-level traffic and emissions monitoring.

